\documentclass[12pt,a4paper]{article}
\usepackage[utf8]{}
\usepackage[T1]{fontenc}


\renewcommand{\baselinestretch}{0.2}

\setlength{\hoffset}{-18pt}         
\setlength{\oddsidemargin}{0pt} % Marge gauche sur pages impaires
\setlength{\evensidemargin}{0pt} % Marge gauche sur pages paires
\setlength{\marginparwidth}{54pt} % Largeur de note dans la marge
\setlength{\textwidth}{481pt} % Largeur de la zone de texte (17cm)
\setlength{\voffset}{-18pt} % Bon pour DOS
\setlength{\marginparsep}{7pt} % Séparation de la marge
\setlength{\topmargin}{0pt} % Pas de marge en haut
\setlength{\headheight}{13pt} % Haut de page
\setlength{\headsep}{4pt} % Entre le haut de page et le texte
\setlength{\footskip}{27pt} % Bas de page + séparation
\setlength{\textheight}{720pt} % Hauteur de la zone de texte (25cm)
\setlength{\parindent}{0pt}

\usepackage{natbib}
\usepackage{graphicx}
\usepackage{amssymb}
\usepackage{amsmath} % Required for the align* environment
\usepackage{float}
\usepackage{pgfplots} 
\usepackage{algorithm}
\usepackage{mathrsfs}

\usepackage{algpseudocode} % test comment
\usepackage[hidelinks]{hyperref}
\pgfplotsset{compat=1.18}


\begin{document}
\author{Yunshan Luo}
\title{DD2380 Artificial Intelligence}
\maketitle
\newpage

\section{Lecture 1a: Introduction to AI}
\textbf{Rational agent approach:} An agent does the thing to maximize its expected performance measure

\vspace{10pt}
\textbf{PEAS} framework:
\begin{itemize}
    \item \textbf{Performane measure:} What is the goal of the agent?
    \item \textbf{Environment: } What is the environment in which the agent operates?
    \item \textbf{Actuators:} What are the actuators of the agent?
    \item \textbf{Sensors:} What are the sensors of the agent?
\end{itemize}

Core definitions:
\begin{itemize}
    \item \textbf{AI (4 viewpoints):} 
    \begin{enumerate}
        \item Thinking humanly
        \item Acting humanly
        \item Thinking rationally
        \item Acting rationally
    \end{enumerate}
    \item \textbf{Agent:} Anything that perceives its environment through sensors and acts through actuators.
    \item \textbf{Agent function:} The mathematical mapting from percept history to action
    \item \textbf{Rational behavior:} Choosing the action that maximizes expected performance measure, given the available information
    \item Perfection rationality is not possible due to imperfect information and computational limits 
\end{itemize}

\vspace{10pt}
\textbf{Gartner's hype cycle:}
\begin{enumerate}
    \item Innovation trigger
    \item Peak of inflated expectations
    \item Trough of disillusionment
    \item Slope of enlightenment
    \item Plateau of productivity
\end{enumerate}

\textbf{Environment properties}
\begin{itemize}
    \item \textbf{Fully vs partially observable:} Does the agent see all relevant states?
    \item \textbf{Determinstic vs stochastic:} Does an action determine the next state? (Stochastic $\implies$ HMMs, MDPs)
    \item \textbf{Episodic vs sequential:} Do earlier actions affect later decisions? (Sequential $\implies$ Planning, RL)
    \item \textbf{Static vs dynamic:} Can the environment change while the agents is deliberating? (Dynamic $\implies$ time pressure, online search)
    \item \textbf{Discrete vs continuous:} Are states/actions/time discrete
    \item \textbf{Single vs multi-agent:} Do other agents' actions matter? (multi-agent $\implies$ game theory/adversarial search)
    \item \textbf{Cooperative vs competitive:} Aligned or opposed objecties? (cooperative $\implies$ multi-agent planning, competitive $\implies$ adversarial search)
\end{itemize}

% continue at lab presentation prep











\end{document}